% Responsables: A, B y C
\section{Conclusiones}

El modelo se estimó con 451 registros únicos y solo seis signos fisiológicos, y
con ellos detectó el 86\,\% de las gestantes de alto riesgo en la base de prueba.
Esto indica que una muestra moderada y un conjunto mínimo de variables medibles
en un puesto rural bastan para sostener una alerta de alto riesgo, aunque no para
graduar los tres niveles. Esto teniendo en cuenta que se tomó la decisión de
eliminar duplicados debido a que al mantenerlos el único efecto era
sobreestimación en las posteriors.

Las estimaciones son confiables desde el punto de vista computacional. Las cuatro
cadenas de Metropolis-Hastings convergieron, con $\hat{R}\leq1.01$ y un ESS mínimo
de 1.003 (glucosa), y las posteriores resultaron mucho más angostas que las
priors. Esto indica que los resultados están determinados por los datos.

En cuanto a las variables, la glucemia fue la de mayor peso ($\mathrm{OR}=3.5$ por
desviación estándar), lo que confirma la hipótesis derivada del análisis
exploratorio, donde ya presentaba la mayor correlación con el riesgo ($r=0.57$).
La presión sistólica ($\mathrm{OR}=1.8$) también confirmó el efecto positivo
esperado a partir de la literatura sobre trastornos hipertensivos. La temperatura
($\mathrm{OR}=2.3$), en cambio, contradijo lo previsto: su prior se centró en cero
por su baja variabilidad, pero ocupó el segundo lugar. Dado que casi todos los
registros se concentran en 98\,°F, este efecto refleja principalmente la presencia
o ausencia de fiebre, se presume que va relacionado con variables que no están
explícitas en el dataset y por lo mismo su explicación puede deberse a relación
con ellas. Además, como se mencionó anteriormente, el modelo es mejor separando
la clase alta y coincide con que los registros con fiebre tienen como etiqueta
riesgo alto.

La regla de decisión determinó buena parte del desempeño. Con la regla MAP, la
clase media casi no se predijo (recall de $0.04$) porque su probabilidad nunca
superó $0.39$, y ajustar un umbral para el alto riesgo tampoco la recuperó: el
recall medio no pasó de $0.08$. La regla de costo esperado elevó el recall de alto
riesgo de $0.68$ a $0.86$, recuperó la clase media hasta $0.40$ y obtuvo la mayor
exactitud balanceada ($0.60$), a cambio de un MCC algo menor ($0.38$ frente a
$0.43$). Esa pérdida proviene del aumento de falsas alarmas, que es el error
definido como menos grave. Se prefirió así un modelo que remite de más a uno que
omite casos graves.

Esta preferencia se justifica por el costo operativo de cada error. Clasificar
como bajo riesgo a una gestante de alto riesgo, el error con mayor peso en la
matriz, implica que no se la remite ni se refuerza su seguimiento, con una
consecuencia potencialmente irreversible. Clasificar como medio a una gestante de
bajo riesgo genera controles adicionales y gastos de transporte, un costo
relevante en un entorno de recursos limitados pero recuperable. Con todo, los
pesos de la matriz son una decisión del grupo y requieren validación clínica.

El bajo desempeño en la clase media sugiere que el supuesto de odds
proporcionales no se cumple del todo: un mismo efecto por variable en ambos
puntos de corte favorece la transición hacia el alto riesgo y deja sin definición
la clase intermedia. Sin embargo, hace parte del trabajo futuro realizar
comparaciones profundas y con otra perspectiva académica.

Además, como la fuente no documenta cómo se asignó RiskLevel, el modelo reproduce
esa etiqueta y no necesariamente un criterio clínico. Si la etiqueta se construyó
con estos mismos signos, parte del desempeño reflejaría la regla de etiquetado.

Por último, la validación externa exige prudencia. Los puntos de corte dependen
de la proporción de clases de esta base, los sensores redondean las mediciones y
la base de prueba es pequeña (117 registros). Aplicar el modelo en otra población
requeriría recalibrarlo y validarlo con datos locales.

El modelo sirve, por tanto, como alerta de alto riesgo que apoya la decisión del
profesional, pero no para graduar el riesgo intermedio. Su mejora depende de una
etiqueta documentada, de modelos que relajen el supuesto de odds proporcionales y
de validación en la población de destino.
